\section{Related Work}


In this section, we review existing works on test-time learning and self-evolving memory. 

\subsection{Test-time Learning}
Test-time learning (TTL) builds upon early work on test-time adaptation (TTA)~\citep{wang2021tent,niu2022efficient,zhang2023memo}, which enables models to adjust to distribution shifts during deployment. Recent advances extend TTA toward \emph{continuous self-improvement}~\citep{iwasawa2021t3a,liu2023tttpp}, allowing models to refine their behavior through online optimization.
Recent \emph{agent-based} studies operationalize such continual improvement via reflection, planning, and self-evolution. Works like \citep{shinn2023reflexion,wang2023voyager,zhou2024eureka,park2023generative,zhao2025selfevolving} and newer frameworks, including \citep{huang2025selfdiscovering,chen2025llmasoptimizer} demonstrate how agents autonomously revise plans, synthesize feedback, and co-evolve \citep{gao2025selfevolvingagents}. These advances mark a shift from static adaptation toward adaptive, self-improving agents capable of continual learning during deployment. Building on this trend, we propose to benchmark such dynamics from a novel \emph{self-evolving memory} perspective.


\subsection{Self-evolving Memory}
Early LLM memory systems primarily served as \emph{passive storage}, maintaining recent dialogues or retrieved facts to compensate for limited context windows~\citep{DBLP:conf/nips/LewisPPPKGKLYR020,DBLP:conf/iclr/AsaiWWSH24,Zhong2023MemoryBankEL,Liu_LlamaIndex_2022,Packer2023MemGPTTL}. Subsequent studies introduced richer management mechanisms, including differentiable read–write controllers~\citep{Modarressi2023RETLLMTA,Liang2023SCMEL} and evaluations under realistic conversational settings~\citep{Maharana2024EvaluatingVL,Wu2024LongMemEvalBC}. Beyond static buffers, recent work explores \emph{policy-driven control}, where the model is explicitly optimized to decide what to store, retrieve, or overwrite~\citep{yu2025memagent,xu2025mem,zhou2025mem1,yan2025memory,Li2025MemOSAO}. Meanwhile, structured memory representations have emerged to organize experiences into relational or procedural forms, as in RepoGraph~\citep{Ouyang2024RepoGraphEA}, MEM0~\citep{mem0}, Zep~\citep{Rasmussen2025ZepAT}, and Dynamic Cheatsheets~\citep{suzgun2025dynamic}. However, there remains no unified evaluation setting and framework for \emph{self-evolving memory}, the ability to reuse and adapt experiences across tasks. \method builds on this trajectory by benchmarking how LLMs not only store and recall but also evolve, reorganize, and reuse memory under streaming task settings.